%==============================================================================
% Curriculum Vitae -- Donghwan Kim
%
% Build:  latexmk -pdf CV_donghwan_kim.tex     (or: make)
% Engine: pdfLaTeX. Only standard CTAN packages, so this also compiles on
%         Overleaf as-is with no extra .cls/.sty files.
%==============================================================================
\documentclass[11pt,letterpaper]{article}

\usepackage[T1]{fontenc}
\usepackage[utf8]{inputenc}
\usepackage{newtxtext,newtxmath}          % Times-like text, matches the old CV
\usepackage[margin=0.8in,top=0.65in,bottom=0.6in]{geometry}
\usepackage{titlesec}
\usepackage{enumitem}
\usepackage{xcolor}
\usepackage[hidelinks]{hyperref}

%------------------------------------------------------------------ formatting
\pagestyle{empty}
\setlength{\parindent}{0pt}
\raggedbottom

% Section headings: small caps, hairline rule underneath.
\titleformat{\section}
  {\large\scshape}{}{0em}{}[\vspace{-0.6em}\rule{\textwidth}{0.6pt}]
\titlespacing{\section}{0pt}{0.8em}{0.35em}

% Tight list used for bullet details under an entry.
\newenvironment{details}
  {\begin{itemize}[leftmargin=1.3em,itemsep=0.12em,topsep=0.2em,parsep=0pt]}
  {\end{itemize}}

% Publication list: the venue tag sits in a fixed-width left column and the
% entry body hangs beside it, so every entry aligns on the same margin.
\newenvironment{publist}
  {\begin{itemize}[leftmargin=5.4em,labelwidth=5.0em,labelsep=0.4em,align=left,
                   itemsep=0.3em,topsep=0.2em,parsep=0pt]%
   \hyphenpenalty=10000 \exhyphenpenalty=10000 \sloppy}
  {\end{itemize}}

%------------------------------------------------------------------- shortcuts
% \entry{left-hand title}{right-hand date}
\newcommand{\entry}[2]{\textbf{#1}\hfill #2\par}
% \pub{TAG} -- starts a publication entry; TAG goes in the left column
\newcommand{\pub}[1]{\item[{[\textbf{#1}]}]}
% own name, bolded in author lists
\newcommand{\me}{\textbf{Donghwan Kim}}
% equal-contribution marker
\newcommand{\eq}{\textsuperscript{*}}
% award annotation
\newcommand{\award}[1]{\textit{(#1)}}
% current month and year, e.g. "October 2026" -- no day, so it never looks stale
\newcommand{\monthyear}{%
  \ifcase\month\or January\or February\or March\or April\or May\or June\or
  July\or August\or September\or October\or November\or December\fi\ \the\year}

\begin{document}

%================================================================= header
\begin{center}
  {\LARGE\bfseries Donghwan Kim}\\[0.45em]
  \href{mailto:djk6434@psu.edu}{djk6434@psu.edu}
  \quad\textbullet\quad \href{https://heyday98.github.io/}{heyday98.github.io}
  \quad\textbullet\quad \href{https://www.linkedin.com/in/dong-hwan-kim/}{linkedin.com/in/dong-hwan-kim}
\end{center}

%================================================================= education
\section{Education}

\entry{The Pennsylvania State University}{State College, PA \quad Aug 2024 -- Present}
Doctor of Philosophy, Computer Science and Engineering \\
Advisor: Kiwan Maeng

\vspace{0.4em}
\entry{Seoul National University (SNU)}{Seoul, Korea \quad Mar 2022 -- Feb 2024}
Master of Science, Interdisciplinary Program in Artificial Intelligence \\
Advisor: Jung Ho Ahn

\vspace{0.4em}
\entry{Seoul National University (SNU)}{Seoul, Korea \quad Mar 2016 -- Feb 2022}
Bachelor of Science, Electrical and Computer Engineering, \textit{cum laude}
\hfill {\footnotesize\itshape Two-year leave for military service}

%================================================================= interests
\section{Research Interests}

My research focuses on building efficient computer systems through hardware--software
co-design. My work spans privacy-preserving machine learning, including differentially
private training and fully homomorphic encryption, and recently extends to LLM
serving systems for agentic AI workloads.

%================================================================= publications
\section{Publications}

\begin{publist}
\pub{IISWC'26} \me, Prakhar Singh, Younghoon Min, Jongryool Kim, Jongse Park,
  Kiwan Maeng, ``Characterizing How Complex Agentic AI Systems Handle General Tasks:
  A Trace-Based Simulation Study,'' \textit{IEEE International Symposium on Workload
  Characterization}, 2026 \award{Best Paper \& Distinguished Artifact Awards}

\pub{OSDI'26} \me, Xin Gu, Jinho Baek, Timothy Lo, Younghoon Min, Kwangsik Shin,
  Jongryool Kim, Jongse Park, Kiwan Maeng, ``Cocoon: A System Architecture for
  Differentially Private Training with Correlated Noises,'' \textit{USENIX Symposium on
  Operating Systems Design and Implementation}, 2026

\pub{CAL'25} SeokHyeon Kong, \me, Euiseong Seo, Kiwan Maeng, ``Characterizing
  the System Overhead of Discrete Noise Generation for Differential Privacy,''
  \textit{IEEE Computer Architecture Letters}, 2025

\pub{CCS'24} Jae Hyung Ju\eq, Jaiyoung Park\eq, Jongmin Kim, Minsik Kang, \me,
  Jung Hee Cheon, Jung Ho Ahn,
  ``NeuJeans: Private Neural Network Inference with Joint Optimization of Convolution
  and FHE Bootstrapping,'' \textit{ACM SIGSAC Conference on Computer and Communications
  Security}, 2024 \award{\eq equal contribution}

\pub{Access'23} \me\eq, Jaiyoung Park\eq, Jongmin Kim, Sangpyo Kim,
  Jung Ho Ahn, ``HyPHEN: A Hybrid Packing Method and Its Optimizations for Homomorphic
  Encryption-based Neural Networks,'' \textit{IEEE Access}, 2023
  \award{\eq equal contribution}

\pub{DISCC'23} Jaiyoung Park, \me, Wonkyung Jung, Sangpyo Kim, Jongmin Kim,
  Jung Hee Cheon, Jung Ho Ahn, ``Toward Practical Privacy-Preserving Convolutional
  Neural Networks Exploiting Fully Homomorphic Encryption,'' \textit{Workshop on Data
  Integrity and Secure Cloud Computing}, 2023

\pub{ISCA'23} Jongmin Kim, Sangpyo Kim, Jaewan Choi, Jaiyoung Park, \me,
  Jung Ho Ahn, ``SHARP: A Short-Word Hierarchical Accelerator for Robust and Practical
  Fully Homomorphic Encryption,'' \textit{International Symposium on Computer
  Architecture}, 2023
\end{publist}

%================================================================= awards
\section{Awards \& Honors}

\entry{Best Paper Award \& Distinguished Artifact Award}{IISWC 2026}

%================================================================= work
\section{Work Experience}

\begin{samepage}
\entry{LLM Serving System Research Intern}{San Jose, CA \quad May 2026 -- Aug 2026}
SK Hynix America
\begin{details}
  \item Built an LLM serving simulator for agentic AI workloads
  \item Explored CXL 3.0 memory pooling for agentic workloads
\end{details}
\end{samepage}

%================================================================= research
\section{Research Experience}

\entry{Graduate Research Assistant}{State College, PA \quad Aug 2024 -- Present}
The Pennsylvania State University \quad (Advisor: Kiwan Maeng)
\begin{details}
  \item \textit{Distributed training systems:} developing a scalable DP-SGD training
    framework on DeepSpeed and PyTorch to optimize multi-GPU gradient aggregation
\end{details}

\vspace{0.5em}
\entry{Graduate Research Assistant}{Seoul, Korea \quad Mar 2022 -- Feb 2024}
Seoul National University \quad (Advisor: Jung Ho Ahn)
\begin{details}
  \item \textit{Algorithm optimization:} devised convolution algorithms and data packing
    formats that minimize expensive homomorphic encryption operations
  \item \textit{AI accelerator design:} designed a domain-specific architecture for
    homomorphic encryption operations
\end{details}

%================================================================= tutorials
\section{Tutorials}

``CrypTorch: Easy-to-use and Efficient PyTorch-based Compiler Infrastructure for
Machine Learning with MPC,'' \textit{ASPLOS 2026 Tutorial}
\href{https://sites.google.com/view/cryptorch-asplos-2026-tutorial}{[link]}

%================================================================= teaching
\section{Teaching Experience}

\entry{Introduction to Computer Architecture (CMPEN 431)}{Fall 2025}
Teaching assistant --- weekly recitations, assignment grading

\vspace{0.5em}
\entry{Computer Organization and Design (CMPEN 331)}{Fall 2024}
Teaching assistant --- weekly recitations and office hours, assignment and exam grading

\vspace{0.5em}
\entry{Programming Methodology}{Spring 2022}
Teaching assistant --- weekly lab sessions, C++ final term project

%================================================================= scholarships
\section{Scholarships}

\entry{AI Fellowship Scholarship}{2022}
Seoul National University, Dept.\ of Engineering

\vspace{0.5em}
\entry{Merit-based Scholarship}{2016, 2017, 2020, 2021}
Seoul National University, Dept.\ of Engineering

\vspace{0.5em}
\entry{Hansung Sonjaehan Scholarship}{2015}
Hansung Sonjaehan Foundation --- \textit{\$4{,}000 for gifted students in science and
engineering}

%================================================================= skills
\section{Technical Skills}

\textit{Machine learning systems:} PyTorch, DeepSpeed, vLLM, Dynamo \\
\textit{Programming languages:} Python, C/C++, CUDA, Verilog \\
\textit{Algorithms \& math:} Differentially private SGD (DP-SGD), Fully Homomorphic
Encryption (CKKS)

\vfill
\begin{center}
  \footnotesize\color{gray} Last updated: \monthyear
\end{center}

\end{document}
